\chapter{Considerações Finais}\label{chap:consideracoes_finais}

Esta proposta parte de uma observação sobre o processo editorial das conferências brasileiras: as convenções que determinam o que a comunidade reconhece como artigo pertinente a uma trilha são tácitas, estão depositadas nos textos já aceitos e raramente se encontram enunciadas nas chamadas de trabalhos. Quem menos domina essas convenções é o pesquisador em início de carreira, e é também quem mais sofre suas consequências. A partir dessa observação, o trabalho propõe um sistema de dois agentes: um que deriva do corpus de artigos aceitos as regras de gênero de uma trilha e verifica a conformidade de um manuscrito a elas, e outro que, a partir de avaliações reais de artigos rejeitados, sinaliza ao autor os riscos de recusa a que seu texto está exposto.

A construção até aqui concluiu os dois fluxos do agente de conformidade de gênero, com a Especificação de Gênero derivada para as três trilhas do corpus. Restam a construção do agente de fronteira de rejeição, a integração das saídas dos dois agentes e a execução do plano de avaliação, previstas para o último trimestre de 2026. A avaliação estabelecerá se cada agente faz o que se propõe, e o fará com dois tipos distintos de evidência: a concordância com anotações de referência e o contraste entre trilhas, no caso do agente de conformidade, e a comparação com decisões editoriais reais, no caso do agente de fronteira.

Duas limitações apontam para o desafio deste trabalho, o corpus do agente de conformidade contém apenas artigos aceitos, de modo que a especificação descreve o gênero tal como realizado pelos textos aprovados, e não a fronteira que separa aceitação de recusa. E o corpus do agente de fronteira, reúne avaliações de conferências internacionais de aprendizado de máquina e processamento de linguagem natural, o que significa que os motivos de recusa que ele reporta refletem essas comunidades, e não uma trilha brasileira. A segunda limitação não decorre de uma escolha de escopo, mas da ausência, hoje, de qualquer conjunto público de avaliações de artigos rejeitados no contexto brasileiro.


\section{Trabalhos futuros}
\label{sec:trabalhos-futuros}

O desdobramento mais relevante é a \textbf{constituição de um corpus de avaliações de artigos rejeitados brasileiros}. Sua ausência é hoje o que mantém o agente de fronteira apoiado em uma base de outro domínio e de outro idioma, e sua existência tornaria os motivos de recusa reportados específicos da comunidade-alvo. Duas vias se apresentam. A primeira é a coleta própria junto a uma conferência da SBC, com consentimento de autores e revisores, aprovação em comitê de ética e observância da legislação de proteção de dados, aproveitando que os pareceres já existem nos sistemas de submissão. A segunda é a identificação de bases públicas de revisão aberta que contemplem artigos rejeitados em domínios mais próximos de Sistemas de Informação do que os hoje disponíveis. Como o grafo de conhecimento é construído no Fluxo A a partir da base de entrada, a substituição consiste em reexecutá-lo sobre o novo corpus.

A \textbf{avaliação de utilidade percebida} com autores em formação constitui o segundo desdobramento. Trata-se de construto perceptual que exige participantes, seus próprios manuscritos e coleta qualitativa, e que só produz percepções interpretáveis depois que a validade interna de cada agente estiver estabelecida. Prevê-se estudo exploratório com um pequeno número de mestrandos, mediante termo de consentimento livre e esclarecido.

A \textbf{extensão da camada retórica ao artigo completo} decorre de uma limitação do modelo adotado. O modelo CARS descreve a organização retórica de introduções, de modo que a camada retórica do artefato restringe-se a essa seção, ao passo que a camada estrutural cobre o documento inteiro. Cobrir método, resultados e discussão exigiria modelos de \textbf{moves} próprios para cada seção, derivados do mesmo corpus.

A \textbf{derivação de especificações a partir de múltiplas edições} permitiria separar o que é regularidade estável de uma trilha do que é próprio de uma edição, como um tema em destaque ou uma composição atípica do comitê. As especificações atuais derivam de uma única edição de cada evento, e séries históricas de anais abertos tornam essa extensão viável sem alteração no método.

A \textbf{validação das regras por revisores experientes} da trilha complementaria os procedimentos previstos. A avaliação planejada atesta que as regras descrevem o corpus e discriminam entre trilhas, não que sejam exaustivas: regras que a comunidade reconhece mas que o corpus não revela permanecem fora do alcance dos procedimentos adotados.

Por fim, a arquitetura admite a \textbf{incorporação de novos agentes} para elevar a qualidade do trabalho.